\section{Introduction}

\label{sec:intro}
%% =========================================================================

Cross-chain bridge protocols face a tension between generality and security. A bridge that supports $n$ chains requires $O(n)$ contract implementations, each a distinct attack surface. Teleport resolves this by separating the protocol into a chain-independent wire format and a chain-dependent adapter layer, connected through a narrow typed interface.

Three prior Lux papers define components used here: Lux Teleport Protocol~\cite{luxtp2022} established lock-and-mint semantics, Lux Bridge~\cite{luxbridge2023} formalized threshold MPC custody, and Lux Warp Messaging~\cite{luxwarp2023} specified intra-ecosystem BLS messaging. This paper unifies these into a single deployable architecture and adds four contributions:

\begin{enumerate}[nosep]
  \item A frozen wire format with exactly four proof tags, audited once and never modified (\S\ref{sec:wire}).
  \item A chain adapter taxonomy with formal adapter interface contracts (\S\ref{sec:adapters}).
  \item A multi-tenant deployment model for white-label bridge instances (\S\ref{sec:multitenant}).
  \item A conformance test framework that enforces correctness at onboarding time (\S\ref{sec:conformance}).
\end{enumerate}

%% =========================================================================
